% Generated draft: requires mathematical and source review.
\documentclass[11pt]{article}
\usepackage[margin=25mm]{geometry}
\usepackage{fontspec}
\setmainfont{Latin Modern Roman}
\usepackage{amsmath,amssymb,mathtools,unicode-math}
\setmathfont{Latin Modern Math}
\usepackage{enumitem,needspace,xurl,hyperref}
\setlength{\parindent}{0pt}
\setlength{\parskip}{5pt}
\setlength{\emergencystretch}{3em}
\newcommand{\N}{\mathbb{N}}
\newcommand{\Z}{\mathbb{Z}}
\newcommand{\Q}{\mathbb{Q}}
\newcommand{\R}{\mathbb{R}}
\newcommand{\C}{\mathbb{C}}
\newcommand{\F}{\mathbb{F}}
\newcommand{\A}{\mathbb{A}}
\newcommand{\PP}{\mathbb{P}}
\newcommand{\E}{\mathbb{E}}
\newcommand{\eps}{\varepsilon}
\newcommand{\dd}{\,\mathrm{d}}
\begin{document}
\section*{UnsolvedMath ID: 30003713}
\textbf{Source problem code:} catalogue entry
\subsection{Definitions and mathematical statement}
Let \(E\) and \(V\) be arbitrary finite-dimensional complex vector spaces. Define the commutative unital complex algebra \(A=\C\oplus E\) by
\[
(\lambda,e)(\mu,f)=(\lambda\mu,\lambda f+\mu e),
\]
so that \(E\) is an ideal with \(E^2=0\). Let \(\mathfrak g=\operatorname{Lie}(V)\) be the free complex Lie algebra generated by \(V\): every complex-linear map from \(V\) to a complex Lie algebra extends uniquely to a Lie algebra homomorphism from \(\mathfrak g\). Form the current Lie algebra
\[
L(E,V)=\mathfrak g\otimes_{\C}A,
\qquad [x\otimes a,y\otimes b]=[x,y]\otimes ab.
\]
Use ordinary Chevalley--Eilenberg homology with trivial coefficients \(\C\). Explicitly, let \(C_n(L)=\bigwedge^n_{\C}L\), with differential
\[
d_n(x_1\wedge\cdots\wedge x_n)
=\sum_{1\leq i<j\leq n}(-1)^{i+j+1}[x_i,x_j]\wedge x_1\wedge\cdots\wedge\widehat{x_i}\wedge\cdots\wedge\widehat{x_j}\wedge\cdots\wedge x_n,
\]
where a hat denotes omission and \(d_0=0\). Write
\[
H_n(L;\C)=\ker(d_n)/\operatorname{im}(d_{n+1}),
\qquad H_*(L;\C)=\bigoplus_{n\geq0}H_n(L;\C).
\]
The coefficient choice is the usual interpretation of the supplied notation \(H_*(L)\); the source record does not explicitly specify coefficients.

\textbf{Problem.} Determine \(H_n(L(E,V);\C)\) for every integer \(n\geq0\) and every pair of finite-dimensional complex vector spaces \(E,V\), preferably by a decomposition into a direct sum of polynomial functors in \(E\) and \(V\), natural under complex-linear maps in both variables. Here a polynomial functor means a functor whose induced maps depend polynomially on the matrix entries of the input linear maps. The source does not specify polynomial degrees, multiplicities, an irreducible decomposition, or whether a finer grading is intended; these remain qualifications of the preferred form of the answer.
\subsection{Short English statement}
For arbitrary finite-dimensional complex vector spaces \(E,V\), compute all homology groups of the free Lie algebra on \(V\) tensored with the square-zero algebra \(\C\oplus E\). Preferably express the answer naturally as a sum of polynomial functors of \(E,V\). Trivial complex coefficients are the conventional interpretation used here; the supplied record leaves them implicit.
\subsection{Sources}
\begin{itemize}
\item Supplied local UnsolvedMath registry record, permanent ID \(30003713\), source problem code \texttt{catalogue entry}, in the Oberwolfach Reports Open Problems collection. The local registry locations are \texttt{data/unsolvedmath/problems.json} and the supplementary \texttt{catalog/problems.json}; the supplied record is the evidence used for this draft.
\item Supplied report citation: \emph{Topology of Arrangements and Representation Stability (2019)}, Oberwolfach Reports, DOI \texttt{10.4171/owr/2018/2}, pp.~43--123. \url{https://doi.org/10.4171/owr/2018/2}. The record describes Oberwolfach Reports as documenting mathematical workshops and their open-problem discussions. It assigns proposed year 2018 and supplies no individual proposer.
\item Supplied literature reference: Zusmanovich, current Lie algebras. \url{https://arxiv.org/abs/0808.0217}.
\item Supplied literature reference: Oberwolfach Report 2/2018. \url{https://doi.org/10.4171/owr/2018/2}.
\item Literature assessment dated exactly \texttt{2026-08-21}: the supplied curation labels the problem open and states that general formulas exist for important low-degree cases of current Lie-algebra homology, but that no general computation of \(H_*(\mathfrak g\otimes A)\) for the stated broad class is available. This is supplied dated literature triage derived from the supplied OWR curation, not an independently verified assessment. Time-sensitive status and cited claims require independent verification.
\item The supplied registry labels the statement \texttt{corrected\_verified} and says that mathematical content was preserved while notation, typography, or omitted definitions were repaired for self-contained publication. This is a supplied editorial qualification, not mathematical verification of this draft.
\end{itemize}
\end{document}
